\chapter{Introduction}

\section{Introduction}
When an earthquake, flood, or fire strikes, the first hours decide how many
people are found alive. Uncrewed aerial vehicles can survey a disaster site far
faster than ground teams, and they do so without risking a rescuer. The imagery
they return, however, is hard to search by eye. A person photographed from
altitude may occupy fewer than ten pixels among rubble, water, or smoke, and a
single sortie can produce thousands of frames. An operator scanning that stream
will miss people. Automating the search is therefore a practical need, not a
convenience, and it is the problem this thesis addresses.

Object detection itself is a mature field. Two-stage detectors in the R-CNN
line~\cite{ren2015fasterrcnn} established high accuracy at high cost, one-stage
detectors in the YOLO family~\cite{redmon2016yolo, jocher2024yolo11} made
detection fast enough for live video, and transformer designs such as
RT-DETR~\cite{zhao2024rtdetr} removed hand-built components. None of this
maturity transfers cleanly to aerial search and rescue. The targets are an
order of magnitude smaller than the objects these detectors were tuned for,
and the published remedies for small objects, extra detection scales, feature
attention, and sliced inference~\cite{akyon2022sahi}, are usually evaluated one
at a time, on generic aerial benchmarks, and without a shared training
protocol. State-space (Mamba) modules~\cite{gu2023mamba} have recently been
promoted for detection as well, again without a controlled test on very small
human targets. Anyone who wants to build a deployable rescue detector today
has to guess which of these additions is worth its cost.

This thesis replaces that guess with a measurement. Starting from the YOLO11m
baseline, it adds one component at a time, attention, a high-resolution
detection scale, and a state-space neck, trains every configuration under one
frozen protocol on the C2A disaster-imagery benchmark~\cite{nihal2024c2a}, and
reports what each change buys in accuracy, parameters, and latency. The outcome
is a compact detector, small enough for airborne hardware, together with an
evidence trail that justifies every design decision in it.

\section{Problem Statement}
Given an aerial image $I$, a detector must output a set of bounding boxes
$\hat{B} = \{(x, y, w, h, c)\}$, one per visible person, each with a confidence
score $c$. A prediction is counted as correct when it overlaps a ground-truth
box $B$ sufficiently, measured by intersection over union,
\begin{equation}
  \mathrm{IoU}(\hat{B}, B) =
    \frac{\mathrm{area}(\hat{B} \cap B)}{\mathrm{area}(\hat{B} \cup B)},
  \label{eq:iou}
\end{equation}
with $\mathrm{IoU} \geq 0.5$ used as the matching threshold throughout this
work. Stated this way the task sounds ordinary. What makes it hard is the
geometry of aerial imagery.

A one-stage detector predicts boxes on a grid whose cell size equals the
feature stride $d$. A target of side $s$ pixels spans
\begin{equation}
  c_d(s) = \frac{s}{d}
  \label{eq:coverage}
\end{equation}
cells in each direction. The finest standard YOLO scale detects at stride
$d = 8$, so a 12-px person spans about 1.5 cells there, and only 0.4 of a cell
at the stride-32 scale. A detector cannot localise what falls inside a
fraction of one cell. In the C2A benchmark used here, about 47\% of the
annotated people are smaller than ten pixels~\cite{nihal2024c2a}, which places
nearly half the dataset at or below the resolving limit of a standard
detection head. Scenes carry 20 to 40 people each, many partially buried or
occluded, against backgrounds of rubble, flood water, fire, and wrecked
vehicles. Figure~\ref{fig:introscene} shows the kind of scene an operator
faces.

\begin{figure}[tbp]
  \centering
  \figorplaceholder{figures/fig_intro_scene.png}{6.5cm}
  \caption{Aerial drone footage of a flood-affected region~\cite{agiindia2024drones}.
  Scenes like this are what a search-and-rescue operator must scan: wide areas,
  cluttered structures, and human figures that occupy only a few pixels each.}
  \label{fig:introscene}
\end{figure}

One more property of the domain shapes the whole thesis. In search and rescue,
a missed person may not get a second chance, while a false alarm costs an
operator a few seconds of review. Recall therefore matters more than
precision, and the operational metric used throughout is the recall-weighted
$F_2$ score,
\begin{equation}
  F_2 = \frac{5PR}{4P + R},
  \label{eq:f2}
\end{equation}
where $P$ is precision and $R$ is recall. The engineering question of this
thesis is then concrete: which architectural changes to a compact one-stage
detector raise recall on sub-ten-pixel people, at a parameter and latency
budget that still fits airborne deployment?

\section{Objectives}
The work pursues the following objectives.
\begin{itemize}
  \item To build a compact detector for very small human targets in aerial
        disaster imagery by refining the YOLO11m baseline through an additive
        ablation, in which channel-and-spatial attention (CBAM), a stride-4
        (P2) detection scale, and a state-space neck are introduced one at a
        time and measured under an identical training and evaluation protocol.
  \item To identify, from those measurements, the configuration that best
        serves very-small-target recall within a deployable parameter and
        latency budget, and to show it performs close to the published state
        of the art on the same benchmark.
  \item To test the recent claim that state-space modules help detection, on
        this task and under this protocol, and to report the outcome whether
        or not it favours the module.
  \item To raise recall on the smallest targets further at inference time,
        without retraining, through slicing-aided inference and test-time
        augmentation.
\end{itemize}

\section{Scope}
The boundaries of the study are set as follows.
\begin{itemize}
  \item \textbf{Dataset.} Training and evaluation use the C2A disaster-imagery
        benchmark~\cite{nihal2024c2a} with its frozen split of 6{,}129
        training, 2{,}043 validation, and 2{,}043 test images, and a single
        object class, \texttt{person}. Cross-dataset generalisation to real
        search-and-rescue collections such as SARD is identified as future
        work, not performed here.
  \item \textbf{Input.} The detector operates on single RGB frames,
        letterbox-resized to $640 \times 640$ for training. Video tracking,
        thermal imagery, and multi-sensor fusion are outside the scope.
  \item \textbf{Architecture.} The study covers attention substitution
        (CBAM), an added stride-4 detection scale, and a state-space neck as
        architectural changes, and SAHI slicing with test-time augmentation as
        inference-time settings. It does not redesign the backbone or search
        architectures automatically.
  \item \textbf{Evaluation.} All results come from the COCO protocol plus an
        operational evaluation at $\mathrm{IoU} = 0.5$, run on a single
        consumer GPU (RTX 4070 Ti SUPER). No field trial with a physical
        drone is included.
\end{itemize}

\section{Unfamiliarity of the Problem}
No prior solution was reused wholesale. The individual ingredients are
published, and are cited where they appear, but the specific design and the
questions asked of it are this work's own.
\begin{itemize}
  \item The additive, single-variable ablation chain on disaster imagery does
        not appear in prior C2A work, which reports whole models against whole
        models~\cite{nihal2024c2a} and leaves open where a gain actually comes
        from.
  \item The state-space experiment places bidirectional local-window Mamba
        blocks inside the \emph{neck} of an otherwise unchanged CNN detector.
        Published Mamba detectors replace the backbone
        instead~\cite{wang2025mambayolo}, which prevents a like-for-like
        attribution.
  \item The attention substitution removes the baseline's C2PSA blocks rather
        than stacking a new module on top, so the attention change arrives at
        a negative parameter cost. Additive attention bolt-ons are the norm in
        the small-object literature.
  \item A negative result is treated as a first-class finding. The state-space
        variant is documented with the same rigour as the successful
        components, because knowing what not to deploy is part of the
        engineering answer.
\end{itemize}

\section{Project Planning}
The thesis spans two academic terms of thirteen weeks each. The first term
covered topic selection, the literature review, preparation of the C2A
dataset, and a benchmark of the YOLO detector family that fixed the baseline.
The second term ran the additive ablation, the state-space exploration, the
inference-time studies, and the analysis and documentation. Writing proceeded
alongside the later experiments rather than after them, so no phase waited on
a finished predecessor. The thesis is the work of a single author, so no
task-assignment matrix across members applies. Figure~\ref{fig:gantt} shows
the timeline.

\begin{figure}[tbp]
  \centering
  \resizebox{\linewidth}{!}{% ============================================================================
%  fig_gantt.tex -- thesis timeline as a pgfgantt chart (Chapter I).
%  Two CSE-4000 terms of 13 weeks each; phases follow the actual project
%  history (family benchmark Jan 2026, additive ablation Feb, state-space
%  exploration Mar-Jun, SAHI/TTA and documentation Jun-Jul).
%  Needs \usepackage{pgfgantt} (loaded in thesisstyle.sty).
%  Include with:  % ============================================================================
%  fig_gantt.tex -- thesis timeline as a pgfgantt chart (Chapter I).
%  Two CSE-4000 terms of 13 weeks each; phases follow the actual project
%  history (family benchmark Jan 2026, additive ablation Feb, state-space
%  exploration Mar-Jun, SAHI/TTA and documentation Jun-Jul).
%  Needs \usepackage{pgfgantt} (loaded in thesisstyle.sty).
%  Include with:  % ============================================================================
%  fig_gantt.tex -- thesis timeline as a pgfgantt chart (Chapter I).
%  Two CSE-4000 terms of 13 weeks each; phases follow the actual project
%  history (family benchmark Jan 2026, additive ablation Feb, state-space
%  exploration Mar-Jun, SAHI/TTA and documentation Jun-Jul).
%  Needs \usepackage{pgfgantt} (loaded in thesisstyle.sty).
%  Include with:  \input{tikz/fig_gantt}
% ============================================================================
\definecolor{gtBlue}{HTML}{6C8EBF}
\definecolor{gtBlueFill}{HTML}{DAE8FC}
\definecolor{gtGreen}{HTML}{82B366}
\definecolor{gtGreenFill}{HTML}{D5E8D4}
\begin{ganttchart}[
    hgrid={*1{draw=black!12}},
    vgrid={*1{draw=black!8}},
    x unit=0.408cm,
    y unit title=0.55cm,
    y unit chart=0.5cm,
    title label font=\scriptsize,
    group label font=\scriptsize\bfseries,
    bar label font=\scriptsize,
    bar height=0.55,
    bar top shift=0.22,
    bar/.append style={fill=gtBlueFill, draw=gtBlue, line width=0.6pt},
    title/.append style={draw=black!30, fill=black!4},
    canvas/.append style={draw=black!30}
  ]{1}{26}
  \gantttitle{First term (weeks)}{13}
  \gantttitle{Second term (weeks)}{13} \\
  \gantttitlelist{1,...,13}{1}
  \gantttitlelist{1,...,13}{1} \\
  \ganttbar{Topic selection and confirmation}{1}{3} \\
  \ganttbar{Literature review}{2}{9} \\
  \ganttbar{Dataset preparation (C2A)}{6}{9} \\
  \ganttbar{Detector family benchmark}{9}{13} \\
  \ganttbar{Pre-defense preparation}{11}{13} \\
  \ganttbar[bar/.append style={fill=gtGreenFill, draw=gtGreen, line width=0.6pt}]
           {Additive ablation (CBAM, P2)}{14}{18} \\
  \ganttbar[bar/.append style={fill=gtGreenFill, draw=gtGreen, line width=0.6pt}]
           {State-space neck exploration}{16}{22} \\
  \ganttbar[bar/.append style={fill=gtGreenFill, draw=gtGreen, line width=0.6pt}]
           {SAHI and TTA evaluation}{20}{24} \\
  \ganttbar{Result analysis and comparison}{21}{25} \\
  \ganttbar{Documentation and defense}{22}{26}
\end{ganttchart}

% ============================================================================
\definecolor{gtBlue}{HTML}{6C8EBF}
\definecolor{gtBlueFill}{HTML}{DAE8FC}
\definecolor{gtGreen}{HTML}{82B366}
\definecolor{gtGreenFill}{HTML}{D5E8D4}
\begin{ganttchart}[
    hgrid={*1{draw=black!12}},
    vgrid={*1{draw=black!8}},
    x unit=0.408cm,
    y unit title=0.55cm,
    y unit chart=0.5cm,
    title label font=\scriptsize,
    group label font=\scriptsize\bfseries,
    bar label font=\scriptsize,
    bar height=0.55,
    bar top shift=0.22,
    bar/.append style={fill=gtBlueFill, draw=gtBlue, line width=0.6pt},
    title/.append style={draw=black!30, fill=black!4},
    canvas/.append style={draw=black!30}
  ]{1}{26}
  \gantttitle{First term (weeks)}{13}
  \gantttitle{Second term (weeks)}{13} \\
  \gantttitlelist{1,...,13}{1}
  \gantttitlelist{1,...,13}{1} \\
  \ganttbar{Topic selection and confirmation}{1}{3} \\
  \ganttbar{Literature review}{2}{9} \\
  \ganttbar{Dataset preparation (C2A)}{6}{9} \\
  \ganttbar{Detector family benchmark}{9}{13} \\
  \ganttbar{Pre-defense preparation}{11}{13} \\
  \ganttbar[bar/.append style={fill=gtGreenFill, draw=gtGreen, line width=0.6pt}]
           {Additive ablation (CBAM, P2)}{14}{18} \\
  \ganttbar[bar/.append style={fill=gtGreenFill, draw=gtGreen, line width=0.6pt}]
           {State-space neck exploration}{16}{22} \\
  \ganttbar[bar/.append style={fill=gtGreenFill, draw=gtGreen, line width=0.6pt}]
           {SAHI and TTA evaluation}{20}{24} \\
  \ganttbar{Result analysis and comparison}{21}{25} \\
  \ganttbar{Documentation and defense}{22}{26}
\end{ganttchart}

% ============================================================================
\definecolor{gtBlue}{HTML}{6C8EBF}
\definecolor{gtBlueFill}{HTML}{DAE8FC}
\definecolor{gtGreen}{HTML}{82B366}
\definecolor{gtGreenFill}{HTML}{D5E8D4}
\begin{ganttchart}[
    hgrid={*1{draw=black!12}},
    vgrid={*1{draw=black!8}},
    x unit=0.408cm,
    y unit title=0.55cm,
    y unit chart=0.5cm,
    title label font=\scriptsize,
    group label font=\scriptsize\bfseries,
    bar label font=\scriptsize,
    bar height=0.55,
    bar top shift=0.22,
    bar/.append style={fill=gtBlueFill, draw=gtBlue, line width=0.6pt},
    title/.append style={draw=black!30, fill=black!4},
    canvas/.append style={draw=black!30}
  ]{1}{26}
  \gantttitle{First term (weeks)}{13}
  \gantttitle{Second term (weeks)}{13} \\
  \gantttitlelist{1,...,13}{1}
  \gantttitlelist{1,...,13}{1} \\
  \ganttbar{Topic selection and confirmation}{1}{3} \\
  \ganttbar{Literature review}{2}{9} \\
  \ganttbar{Dataset preparation (C2A)}{6}{9} \\
  \ganttbar{Detector family benchmark}{9}{13} \\
  \ganttbar{Pre-defense preparation}{11}{13} \\
  \ganttbar[bar/.append style={fill=gtGreenFill, draw=gtGreen, line width=0.6pt}]
           {Additive ablation (CBAM, P2)}{14}{18} \\
  \ganttbar[bar/.append style={fill=gtGreenFill, draw=gtGreen, line width=0.6pt}]
           {State-space neck exploration}{16}{22} \\
  \ganttbar[bar/.append style={fill=gtGreenFill, draw=gtGreen, line width=0.6pt}]
           {SAHI and TTA evaluation}{20}{24} \\
  \ganttbar{Result analysis and comparison}{21}{25} \\
  \ganttbar{Documentation and defense}{22}{26}
\end{ganttchart}
}
  \caption{Timeline of the thesis across two thirteen-week terms. First-term
  phases (blue) built the foundation and the baseline benchmark; second-term
  phases (green and blue) ran the ablation, the state-space exploration, the
  inference-time studies, and the documentation.}
  \label{fig:gantt}
\end{figure}

\section{Contribution}
The main contributions of this thesis are the following.
\begin{itemize}
  \item A four-configuration additive ablation on C2A, trained and evaluated
        under one frozen protocol, that attributes every gain to a single
        architectural change.
  \item A CBAM substitution that replaces the baseline's C2PSA attention and
        removes roughly one million parameters while matching or improving
        the operational scores, with the lowest latency of the four
        configurations.
  \item A stride-4 (P2) detection scale whose effect is quantified rather
        than assumed: $\mathrm{AP}_{50}$ rises from 0.843 to 0.853 and
        very-tiny recall from 0.743 to 0.757, at about half a million
        parameters and one millisecond of latency.
  \item A controlled negative result for the state-space neck: 2.4 million
        added parameters and 2.8 times the latency for no accuracy change,
        which argues against adopting such modules for this task on current
        evidence.
  \item An inference-time study showing that test-time augmentation at
        1280-px input lifts very-tiny recall from 0.758 to 0.850 at 60~ms per
        image, a larger gain than any slicing configuration at a fraction of
        the slicing cost.
  \item A recommended detector of 19.6 million parameters that matches or
        exceeds every published C2A baseline except YOLOv9-e while using
        about a third of that model's parameters.
\end{itemize}

\section{Application of the Work}
The detector developed here has direct uses beyond the motivating scenario.
\begin{enumerate}
  \item \textbf{Search and rescue triage.} The immediate application: a
        drone or a laptop at an incident site flags likely victims in live
        aerial imagery so operators review candidates instead of raw frames.
  \item \textbf{Disaster damage assessment.} The same detections, aggregated
        over a survey flight, tell relief planners where people are
        concentrated after a flood or earthquake.
  \item \textbf{Crowd monitoring.} Counting very small people from altitude
        is the same technical problem, so the model transfers to public-event
        safety monitoring.
  \item \textbf{Wide-area security.} Border and coastline patrols scan large
        areas for isolated people, a low-density variant of the same task.
  \item \textbf{Embedded deployment.} At 19.6 million parameters the design
        fits consumer and edge GPUs, and the same recipe transfers to a
        smaller YOLO11 backbone for fully on-board operation.
\end{enumerate}

\section{Organization of Report}
The remainder of the thesis is organised as follows.
\begin{itemize}
  \item \textbf{Chapter II} reviews the detector families, the small-object
        and aerial detection literature, and the base architectures this work
        builds on, and closes with the research gap.
  \item \textbf{Chapter III} presents the proposed detector: the system
        overview, the layer-level architecture, the rationale for each
        novelty, the loss design, and the inference-time procedures.
  \item \textbf{Chapter IV} reports the experimental setup, the additive
        ablation results, the per-size analysis, the comparison with
        published detectors, and the SAHI and TTA studies.
  \item \textbf{Chapter V} discusses the societal, ethical, legal, and safety
        aspects of deploying a person detector.
  \item \textbf{Chapter VI} maps the work onto the complex engineering
        problem and activity attributes it satisfies.
  \item \textbf{Chapter VII} summarises the findings, states the limitations,
        and sets out future work.
\end{itemize}
